\SourceLeaf{25}
\par
Die Anziehung in der zur $y$-Koordinatenachse parallelen und ihr entgegengesetzten Richtung ergibt sich, wenn man in diesen Formeln $a$ durch $b$ ersetzt. Diese beiden Kräfte sind demnach einer einzigen Kraft gleichwertig, deren Richtung senkrecht zur Achse $2C$ ist und deren Stärke man erhält, wenn man in der eben angegebenen Formel $a$ durch den Abstand des angezogenen Punkts von dieser Achse ersetzt.

Schließlich ergibt sich die Anziehung in der zur $z$-Koordinatenachse parallelen und ihr entgegengesetzten Richtung, also senkrecht zum Äquator, im Fall $B=A$ zu
\[
=\frac{4\pi cAA}{CC}\int\frac{tt\mathrm{d}t}{1-(1-\frac{AA}{CC})tt}
\]
woraus sich, falls $C<A$, mit der obigen Setzung $\frac{C}{A}=\cos\varphi$ ergibt
\[
=\frac{4\pi c\cos\varphi}{\sin\varphi^3}(\operatorname{tang}\varphi-\varphi)
\]
falls aber $C>A$, ergibt sich
\[
\frac{4\pi cAAC}{(CC-AA)^{\frac32}}\log\frac{C+\BareRoot{(CC-AA)}}{A}-\frac{4\pi cAA}{CC-AA}
\]
Sind zuletzt alle drei $A,B,C$ gleich, ist der Körper also eine Kugel, so ergeben sich für die Anziehungen in den drei Hauptrichtungen
\[
\tfrac43\pi a,\quad\tfrac43\pi b,\quad\tfrac43\pi c
\]
Sie stimmen also mit den Anziehungen überein, die ein kugelförmiger Kern, auf dessen Oberfläche der angezogene Punkt liegt, ausüben würde, wenn seine Masse im Mittelpunkt konzentriert wäre. Daraus folgt unmittelbar, dass \textit{äußere} Punkte von einer Kugel so angezogen werden, als läge ihre gesamte Masse im Mittelpunkt, wie \textsc{Newton} zuerst gezeigt hat.

\par\bigskip\Needspace{6\baselineskip}
\pdfbookmark[1]{Zusatz}{attractionis-additamentum}
\centerline{ZUSATZ.}\medskip
\par
Nachdem dies bereits niedergeschrieben war, wurde durch einen Hinweis des bedeutenden \textsc{Laplace} eine hervorragende Abhandlung des angesehenen \textsc{Ivory} in den \textit{Philosophical Transactions} für 1809 bekannt. Darin wird derselbe Gegenstand nach einer Methode behandelt, die sich von denen des bedeutenden \textsc{Laplace} und \textsc{Legendre} völlig unterscheidet. Mit größter Eleganz zeigte jener Geometer, wie die Anziehung auf einen äußeren Punkt auf die Anziehung auf einen inneren Punkt zurückgeführt werden kann, wie sich also der stets für schwieriger gehaltene Teil der Aufgabe auf den leichteren zurückführen lässt. Die Methode aber, nach der
